% Faithful transcription — original Italian (18th c.). Transcribed from the 1911 Opere
% Matematiche, vol. II (Wikimedia Commons djvu), the complete memoir "Metodo per misurare la
% lemniscata, Schediasma I", printed pages 293–297 (= djvu pages 323–327). \origpage uses the
% PRINTED page numbers. Content and archaic spelling (ànno, à, Bernulli, elisse, abscisse,
% elittico) follow the print; only markup is standardized (PLAN.md §4.4). NOTE: eq. (2) is
% reproduced faithfully with the denominator \sqrt{a^2-az} exactly as printed (it differs from
% the \sqrt{a^2-z^2} elsewhere — a probable misprint left for editorial review), and Fagnano's
% era notation "zz" for z^2 is kept where the print uses it.
\documentclass{article}
\usepackage{readmasters}
\begin{document}

\origpage{293}
\section*{XXXII. Metodo per misurare la lemniscata}
\subsection*{Schediasma I (*)}
\emph{(*) Giornale de' letterati d'Italia, tomo XXIX, pag.~258.}

I due sommi geometri signor Giacomo e signor Giovanni fratelli Bernulli ànno renduta
celebre la lemniscata, servendosi de' suoi archi per costruire l'isocrona paracentrica,
come può vedersi negli Atti di Lipsia dell'anno 1694. Egli è visibile, che misurando la
lemniscata mediante qualche altra curva di lei più semplice, si ottiene una costruzione
più perfetta non solo dell'isocrona paracentrica, ma ancora delle altre infinite curve,
che per essere costruite possono dipendere dalla lemniscata; e però mi lusingo, che non
sieno per dispiacere agl'intendenti le misure di questa curva da me scoperte, le quali
esporrò successivamente in due schediasmi.

\emph{Supposizioni note agl'intendenti del calcolo infinitesimale.} --- Sia la lemniscata
$CQACFC$ (fig.~24), il cui semiasse $CA = a$; si sa, che prendendo nel centro $C$ l'origine
dell'abscisse $(x)$, e chiamando $(y)$ le ordinate che sono all'asse normali, la natura
della lemniscata s'esprime con quest'equazione $x^{2} + y^{2} = a\sqrt{x^{2} - y^{2}}$.

Si sa ancora che se si chiama $z$ la corda indeterminata $CQ = \sqrt{x^{2} + y^{2}}$, si à
l'arco diretto $CQ = \displaystyle\int \frac{a^{2}\,dz}{\sqrt{a^{4} - z^{4}}}$, e l'arco inverso
\[ QA = \mathrm{arc.}\,CA - \mathrm{arc.}\,CQ = \int \frac{a^{2}\,dz}{\sqrt{a^{4} - z^{4}}}. \]

Sia l'elisse $ADFNA$ (fig.~24), il cui semiasse minore $CF = a$, e il semiasse maggiore
$CD = a\sqrt{2}$; chiamasi $z$ l'abscissa indeterminata $CH$, che ha l'origine del centro
$C$ dell'elisse, ed è uguale alla corda $CQ$ della lemniscata, e tirisi l'ordinata $HI$
parallela all'asse maggiore; è già noto che l'arco diretto $DI$ di quest'elisse ha per sua
espressione
\[ \int dz\, \frac{\sqrt{a^{2} + z^{2}}}{\sqrt{a^{2} - z^{2}}}, \]
\origpage{294}
e che l'arco inverso
\[ IF = \mathrm{arc.}\,DF - \mathrm{arc.}\,DI = \int -dz\, \frac{\sqrt{a^{2} + z^{2}}}{\sqrt{a^{2} - z^{2}}}. \]

\rmfigure{figures/fig-24.png}{Fig.~24.}{Fig. 24: the lemniscate CQACFC and its circumscribed ellipse ADFNA, from the original plate (Tav. II, Opere Matematiche vol. II).}

Sia finalmente l'iperbola equilatera $LMP$ (fig.~25), il cui semiasse $SM = a$. Si nomini
$t$ l'applicata indeterminata $SO$, che partendo dal centro $S$ tagli l'arco $MO$; si sa,
che questo medesimo arco s'esprime così
\[ \int \frac{t^{2}\,dt}{\sqrt{t^{4} - a^{4}}}. \]

\rmfigure{figures/fig-25.png}{Fig.~25.}{Fig. 25: the equilateral hyperbola LMP with ordinate SO and arc MO, from the original plate (Tav. III, Opere Matematiche vol. II).}

\textbf{Teorema I.} --- Sieno le due infrascritte equazioni (1), e (2). Io dico, che posta
la prima di esse, sussiste anche l'altra

\[ t = a\, \frac{\sqrt{a^{2} + z^{2}}}{\sqrt{a^{2} - z^{2}}}. \tag{1} \]

\[ \int \frac{a^{2}\,dz}{\sqrt{a^{4} - z^{4}}} = \int dz\, \frac{\sqrt{a^{2} + zz}}{\sqrt{a^{2} - az}} + \int \frac{t^{2}\,dt}{\sqrt{t^{4} - a^{4}}} - \frac{zt}{a}. \tag{2} \]

La verità di questo teorema si dimostra differenziando, e sostituendo in luogo di $t$, e di
$dt$ i loro valori in $z$, e $dz$ tratti dall'equazione (1).

\textbf{Corollario} (fig.~24 e 25). --- Se nella lemniscata la corda $CQ = z$, e nell'elisse
l'abscissa $CH$ è anche essa $= z$, e nell'iperbola equilatera $LMP$ l'applicata centrale
$SO = t$; assegnando a $t$ il suo valore espresso nell'equazione (1) e ponendo nell'equazione
(2) gli archi delle curve invece delle loro espressioni già dichiarate nelle supposizioni
superiori, s'ottiene

\[ \mathrm{Arc.}\,CQ = \mathrm{arc.}\,DI + \mathrm{arc.}\,MO = \frac{zt}{a}. \tag{3} \]

\textbf{Scolio I.} --- Supponendo $z = 0$, gli archi diretti $CQ$ della lemniscata, e $DI$
dell'elisse sono uguali a zero; è nullo ancora l'arco iperbolico $MO$, poichè in virtù
dell'equazione (1) l'annullamento di $z$ rende $SO\,(t) = a = SM$; s'annienta eziandio
l'espressione rettilinea, e tutto ciò è un certo indizio, che l'equazione (3) è completa, e
non abbisogna dell'aggiunta, o sottrazione d'alcuna quantità costante. Ma facendo $z = a$,
allora l'arco diretto della lemniscata diventa eguale all'arco $CQA$, e l'arco elittico $DI$
all'arco $DIF$; ma per l'equazione (1) l'applicata centrale $SO\,(t)$ dell'iperbola diviene
in questo caso infinita, rendendo infinito anche l'arco iperbolico $MO$, e però sembra
impossibile di giungere a conoscere per questa strada il
\origpage{295}
valore del quadrante della lemniscata espresso in quantità finite, e per conseguenza la
misura di questa curva non può dirsi ancora perfettamente scoperta.

\textbf{Teorema II.} --- Sieno le due equazioni infrascritte (4) e (5); io dico, che posta
la prima di esse, sussiste anche l'altra

\[ r = \frac{a^{2}}{z}. \tag{4} \]

\[ \int -\frac{a^{2}\,dz}{\sqrt{a^{4} - z^{4}}} = \int -dz\, \frac{\sqrt{a^{2} + z^{2}}}{\sqrt{a^{2} - z^{2}}} + \int \frac{r^{2}\,dr}{\sqrt{r^{4} - a^{4}}} - \frac{1}{z}\sqrt{a^{4} - z^{4}}. \tag{5} \]

La differenziazione e la sostituzione de' valori di $r$, e $dr$ in $z$, e $dz$ dedotti
dall'equazione (4) dimostrano questo teorema.

\textbf{Corollario I} (fig.~24 e 25). --- Facciasi, come sopra, nella lemniscata la corda
$CQ = z$, nell'elisse l'abscissa $CH$ anch'essa $= z$, e nell'iperbola l'applicata centrale
$SL = r$, attribuendo ad $r$ il suo valore notato nell'equazione (4), e surrogando
nell'equazione (5) gli archi delle curve in vece delle loro espressioni, conforme si è fatto
nel corollario del I teorema, ritrovasi

\[ \mathrm{Arc.}\,QA = \mathrm{arc.}\,IF + \mathrm{arc.}\,ML - \frac{1}{z}\sqrt{a^{4} - z^{4}}. \tag{6} \]

\textbf{Scolio II.} --- La supposizione di $z = a$ rende nulli gli archi inversi $QA$ della
lemniscata, ed $IF$ dell'elisse; annulla ancora l'espressione rettilinea, e l'arco iperbolico
$ML$, atteso che in vigore dell'equazione (4) l'applicata centrale dell'iperbola $SL\,(r)$
diviene uguale al semiasse $SM\,(a)$, dunque l'equazione (5) è completa. Ma la supposizione
di $z = 0$, prolunga in infinito l'applicata centrale dell'iperbola, e l'arco, di essa, nello
stesso tempo, che fa divenire gli archi inversi della lemniscata, e dell'elisse uguali
rispettivamente ai quadranti $ACQ$, $FID$; ed ecco rinascere il medesimo inconveniente
esposto nel fine del I scolio; ma questa difficoltà resterà superata dal corollario, che
segue. Intanto si noti; primo, che tirando nell'iperbola la $OP$ parallela al secondo asse,
l'arco $LMP$ è uguale ai due archi $MO$, $ML$. Secondo, che assegnando a $t$ il suo valore
espresso nell'equazione (1), questa quantità complessa $\dfrac{zt}{a} + \dfrac{1}{z}\sqrt{a^{4} - z^{4}}$
equivale a quest'altra quantità più semplice $\dfrac{at}{z}$, come si esperimenterà col calcolo.
\origpage{296}
\textbf{Corollario II} (fig.~24 e 25). --- Aggiungendo le due equazioni (3) e (6), e
servendosi dei due avvertimenti notati nel fine del precedente scolio, si scuopre
$\mathrm{arc.}\,CQA = \mathrm{arc.}\,DIF + \mathrm{arc.}\,LMP - \dfrac{at}{z}$; e finalmente
moltiplicando per 4 quest'ultima equazione, ritrovasi che l'intera periferia della lemniscata
è uguale all'intera periferia dell'elisse circoscritta $ADFNA$, più il quadruplo dell'arco
$LMP$ dell'iperbola equilatera meno $\dfrac{4at}{z}$ ch'è una nuova, ed egregia proprietà di
queste curve.

\textbf{Teorema III.} --- Sieno le due equazioni infrascritte (7) e (8); io dico, che posta
la prima di esse, sussiste anche l'altra

\[ u = a\, \frac{\sqrt{a^{2} - z^{2}}}{\sqrt{a^{2} + z^{2}}}. \tag{7} \]

\[ \int \frac{a^{2}\,dz}{\sqrt{a^{4} - z^{4}}} = \int -\frac{a^{2}\,du}{\sqrt{a^{4} - u^{4}}}. \tag{8} \]

Questo teorema si dimostra, come gli altri due, che lo precedono, differenziando, e facendo
le debite sostituzioni tratte dall'equazione (7).

\textbf{Corollario I} (fig.~24). --- Prendasi nella lemniscata la corda $CQ = z$, e l'altra
corda $CE = u$, attribuiscasi ad $u$ il suo valore espresso nell'equazione (7), si
sostituiscano nell'equazione (8) gli archi diretto, ed inverso della curva in luogo delle
loro espressioni, e si troverà $\mathrm{arc.}\,CQ = \mathrm{arc.}\,EA$.

\textbf{Scolio III} (fig.~24). --- Se $z = 0$ l'arco diretto $CQ$ è nullo, e in virtù
dell'equazione (7) l'annullamento di $z$ rende $CE\,(u) = a$, di modo che l'arco inverso $EA$
è nullo anch'esso; il che fa conoscere, che l'equazione notata nel fine del precedente
corollario è completa. Questa insigne proprietà della lemniscata è stata da me scoperta, ed
esposta in maniera diversa nel mio metodo per rettificare la differenza degli archi in infinite
specie di parabole irrettificabili.

\textbf{Corollario II} (fig.~24). --- Trovando, come qui sopra ho insegnato di fare, l'arco
inverso $EA$ eguale all'arco diretto $CQ$, potrà misurarsi questo medesimo arco $CQ$ mediante
il I.\ corollario del II teorema, e viceversa trovando l'arco diretto $CQ$ eguale all'arco
inverso $EA$, potrà essere misurato questo medesimo arco $EA$ per mezzo del corollario del I
teorema.
\origpage{297}
\textbf{Corollario III.} --- \textbf{Problema.} --- Tagliare per mezzo il quadrante della
lemniscata (fig.~24).

Se la corda $CQ\,(z)$ della lemniscata è uguale all'altra corda $CE\,(u)$ di essa, e se nel
tempo stesso $u$ ha il suo valore notato nell'equazione (7), allora
$z = u = a\, \dfrac{\sqrt{a^{2} - z^{2}}}{\sqrt{a^{2} + z^{2}}}$, donde risulta fatte le debite
operazioni
\[ z = u = \sqrt{a^{2}\sqrt{2} - a^{2}} \]
e però attribuendo a $z$ quest'ultimo valore, i due punti della lemniscata $Q$ ed $E$
coincideranno in $B$, in maniera che l'arco diretto $CB$ sarà eguale all'arco inverso $BA$, e
la corda $CB = u = z$ taglierà per mezzo il quadrante della lemniscata. Il che dovea
ritrovarsi.

\end{document}
